\chapter{From the page to a distribution over sounds}

\textbf{Purpose:} Follow one Salishan word from the mark on the page to a symbol a distribution can
be taken over, and say which stage each tool in this tree belongs to, so that a reader measuring
anything on this corpus knows what their symbols stand for.
\textbf{Scope:}
\texttt{maint/\allowbreak{}data/\allowbreak{}salishan/\allowbreak{}corpus\_\allowbreak{}script\_\allowbreak{}extraction/\allowbreak{}salish\_\allowbreak{}marking.\allowbreak{}py},
\texttt{maint/\allowbreak{}data/\allowbreak{}salishan/\allowbreak{}corpus\_\allowbreak{}script\_\allowbreak{}extraction/\allowbreak{}paper\_\allowbreak{}config.\allowbreak{}py},
\texttt{build/\allowbreak{}experiments/\allowbreak{}salish\_\allowbreak{}phonemes.\allowbreak{}py},
the hand extractions under \texttt{oracles/\allowbreak{}} in the corpus.

This chapter exists because the question keeps being asked and answered in conversation. The answer
lives in a header comment in \texttt{build/\allowbreak{}experiments/\allowbreak{}}, which is not
where anybody looks, and the reader who needs it is the reader who does not yet know it is there.

\section{Three stages, and they are not interchangeable}

A symbol in this corpus can mean three different things, and a distribution taken over each of them
answers a different question.

\begin{enumerate}
  \item A \textbf{character} is a codepoint in a file. A distribution over characters is a
        distribution over what a transcriber typed.
  \item A \textbf{segment} is a base with everything that modifies it. A distribution over segments
        is a distribution over the units the orthography writes.
  \item A \textbf{sound} is what the segment stands for, with the spellings of one sound collapsed.
        A distribution over sounds is the first one that is about the language.
\end{enumerate}

Each stage discards something the stage before it carried, and each discard is a claim that has to
be true. Skipping a stage does not leave the measurement neutral. It leaves it measuring the thing
that stage was there to remove.

\section{Stage zero, the page}

Nothing above starts until somebody has read the paper. A text extraction of a scanned typescript is
not the paper: all four ICSNL 2 papers held here carry none of their own orthography in their text
layer, zero of the nine marks the corpus uses, and
\texttt{ORTHOGRAPHY\_\allowbreak{}ABSENT} in
\texttt{maint/\allowbreak{}data/\allowbreak{}salishan/\allowbreak{}hand\_\allowbreak{}extraction/\allowbreak{}papers.\allowbreak{}py}
names them so that no check grades a correct reading against a file the reading is right to differ
from.

The trap runs the other way as well. \emph{The Moon and the Birchbark Canoe} is born digital and its
text layer is faithful, but the embedded font draws U+026C, the l with a belt, as an l with a bar
through the stem. A reader working from page renders sees ł and writes ł. The paper's text layer
holds U+026C 112 times and U+0142 not once. The glyph on the page and the character the page is
drawing are two different questions, and the hand extraction records the second with the first
written down beside it.

\section{Stage one, characters}

\texttt{TEXT\_\allowbreak{}SPACE} is the union of every character these orthographies write their
languages with
(\texttt{salish\_\allowbreak{}marking.\allowbreak{}py:88}). It is built from
\texttt{MARKED}, the marked consonants and the schwa
(\texttt{salish\_\allowbreak{}marking.\allowbreak{}py:30}), \texttt{PRACTICAL}, the digit 7
(\texttt{salish\_\allowbreak{}marking.\allowbreak{}py:76}), and the stacked glottalization marks,
the labialization modifier and Nater's raised space.

Two things about it are worth stating plainly, because both have caused a wrong measurement here.

\textbf{It is a membership test and not an alphabet.} Its job is to answer whether a token is in one
of these languages. A per-paper set is written as this plus what that paper adds, never as that
paper's own alphabet, and the reason is in the file: a test built on one paper's characters reads
every other paper as holding no language, and then reports nothing missing from them either.

\textbf{It carries no phonetic content at all.} \texttt{ɬ} and \texttt{ł} are both in it, as separate
members, because both occur. Nothing in \texttt{TEXT\_\allowbreak{}SPACE} says they are one sound.

Typographic ligatures are repaired before anything else, without a test, because \texttt{fi} is
\texttt{fi} and no decision is involved
(\texttt{salish\_\allowbreak{}marking.\allowbreak{}py:56}). That one was found by a case sensitive
match turning up five against five.

\section{Stage two, segments}

\texttt{ʷ} is not a segment. It is labialization on the consonant before it. \texttt{̓} is not a
segment either; it is glottalization on the base it sits over. Counting per character measures two
halves of one segment as though they were independent.

The rule is to join every combining mark and every modifier letter to the base in front of it. The
hand extractions already do this when they write a form. Applied to Bev Phillips's
\emph{When Old One Created the Earth}, it turns 5552 characters over 35 distinct into 4839 segments
over 50 distinct. The inventory grows under clustering instead of shrinking, because \texttt{kʷ} and
\texttt{k} stop being one letter counted twice. Twenty of the fifty carry a mark or a modifier:

\begin{quote}
\texttt{c̓ kʷ k̓ k̓ʷ m̓ n̓ p̓ qʷ q̓ q̓ʷ w̓ xʷ x̣ x̣ʷ y̓ z̓ ƛ̓ ə́ ʕʷ ʕ̓ʷ}
\end{quote}

\texttt{ə́} is one of the fifty. It is the segment that carried the Lushootseed dialect border. At
the character level it is two codepoints, and the radix found it at width two for that reason.

There is no shared function in this tree that performs this clustering. The rule is stated in prose,
followed by hand in the oracles, and implemented once inside
\texttt{build/\allowbreak{}experiments/\allowbreak{}pooled\_\allowbreak{}sound.\allowbreak{}py}.
Anyone measuring segments today writes it again.

\section{Stage three, sounds}

\texttt{build/\allowbreak{}experiments/\allowbreak{}salish\_\allowbreak{}phonemes.\allowbreak{}py}
holds the table that takes a spelling and returns the sound it stands for. Every asking of this
question arrives at that file, and nothing in the tree points to it.

The spellings that differ across the corpus, counted by segment over the hand extractions that
existed when the table was built, which was 23 of them:

\tiny
\setlength{\tabcolsep}{2pt}
\begin{longtable}{@{}>{\arraybackslash}p{\dimexpr\textwidth/4-2\tabcolsep\relax}>{\arraybackslash}p{\dimexpr\textwidth*3/4-2\tabcolsep\relax}@{}}
\toprule
Sound & How the corpus writes it \\
\midrule
lateral fricative & \texttt{ɬ} in seventeen papers, \texttt{ł} in four. LaFontaine and Janzen write \texttt{ł} 329 times and \texttt{ɬ} never. Nater's voiceless paper writes \texttt{ł} 98 times and \texttt{ɬ} three. \\
glottal stop & \texttt{ʔ} in most, the digit \texttt{7} in the van Eijk orthography. Alexander and Davis write \texttt{7} 934 times and \texttt{ʔ} four. Matthewson and Redan write \texttt{7} 125 times and \texttt{ʔ} twice. A capital \texttt{Ɂ} occurs in one paper. \\
uvular fricative & \texttt{χ}, \texttt{x̌} and \texttt{x̣}, all three in six or more papers. \\
schwa & \texttt{ə} at U+0259 and \texttt{ǝ} at U+01DD, a turned e and a different codepoint. \\
\bottomrule
\end{longtable}
\normalsize

Three of those four were named in the refs chapter. The fourth came out of counting, and so did the
capital glottal stop and the aspirated \texttt{pʰ} and \texttt{tʰ} that only Nater writes. That is
the argument for building the table by counting the corpus instead of from a description of it.

\subsection{Where the correspondences come from}

Nater states his own, in a paper this corpus holds: \texttt{/ł/ = /ɬ/}, \texttt{/χ/ = /x̌/},
\texttt{/χʷ/ = /x̌ʷ/}, \texttt{/p’, t’, q’/ = /p̓, t̓, q̓/}, \texttt{/kʷ’, qʷ’/ = /k̓ʷ, q̓ʷ/} and
\texttt{/’/ = [ʔ]}. A linguist writing down what his own symbols equal is the strongest source
available here. The van Eijk 7 is declared in the tree already. The rest follow the standard
descriptions of Americanist transcription, they were assigned by hand, and any error in them is an
error in every result taken downstream.

\subsection{Two errors that table had, both caught by reading its own output}

The first draft sent \texttt{kʷ} and \texttt{qʷ} to one symbol. A labialized velar and a labialized
uvular are different phonemes in these languages, and the collapse showed itself at once:
\texttt{kʷɬ}, \texttt{qɬ} and \texttt{qʷł} arrived as one form, which merges three words.

The second draft gave the uvular fricative the symbol \texttt{x}. A plain velar \texttt{x} already
passes through as that same symbol. The mapped uvular and the unmapped velar then met in one cell.
A pass-through is not free of the namespace it passes into.

Every sound now takes a codepoint of its own from the private use area, assigned by the table and
never typed, so no sound can collide with a letter the corpus writes. Both errors are recorded
instead of quietly fixed, because a table like this is exactly where a wrong assertion is invisible:
nothing downstream can tell a real merger from a typo.

\subsection{What the mapping does not claim}

Mapping \texttt{ɬ} and \texttt{ł} to one symbol says two spellings are one sound. That is a fact
about transcription practice. Whether Nuxalk's lateral fricative and Lushootseed's are the same
sound is a question for a comparativist, and nothing here asserts it. What the mapping buys is
narrow and it is enough: a difference measured afterward is not the difference between two typists.

\section{What a distribution reads at each stage}

\tiny
\setlength{\tabcolsep}{2pt}
\begin{longtable}{@{}>{\arraybackslash}p{\dimexpr\textwidth/4-2\tabcolsep\relax}>{\arraybackslash}p{\dimexpr\textwidth*3/4-2\tabcolsep\relax}@{}}
\toprule
Over what & What the number is about \\
\midrule
bytes & the encoding, and whether a mark is precomposed or decomposed \\
characters & the transcriber, and which conference volume the paper appeared in \\
segments & the orthography, comparable inside one paper and across papers sharing a convention \\
sounds & the language, comparable across every paper in the corpus \\
\bottomrule
\end{longtable}
\normalsize

The failure this ordering exists to prevent has a shape worth naming. A model counting character
bigrams reads two dialects that share a word as maximally far apart, because it is measuring the
transcriber and not the speaker. The distance is real, it is stable, it will replicate, and it is
about the wrong thing.

\section{What is not built}

Three gaps, stated so that nobody has to rediscover them.

\textbf{The phoneme table lives in an experiment.} It is under
\texttt{build/\allowbreak{}experiments/\allowbreak{}} and nothing in the engine imports it. Work
needing sounds today either finds that file or writes the table again.

\textbf{The clustering rule has no shared implementation.} Stage two is prose plus one copy inside
one experiment.

\textbf{The word web is still built on written forms.} The refs chapter states this and it is
unchanged: nothing takes a character out of \texttt{TEXT\_\allowbreak{}SPACE} and hands back what it
stands for at the point the web is built, so two dialects writing one word two ways are still two
nodes. Stage three is what would close it.

\textbf{Author:} dstroy0 (Douglas Quigg) <dquigg123@gmail.com>
\textbf{Date:} 2026-09-10
